\section{Instrumenting the measurement campaign: observability of the
platform itself}
\label{sec:campagne}

Every result in Sections~\ref{sec:ovstarvation} and~\ref{sec:sqrtvinsint} rests on
a premise that was never itself measured: that during the run which produced
the numbers, the platform was delivering data at the rate the estimators
assume. This section closes that gap. It derives the thresholds of a
health monitor from the statistics of the streams rather than from round
numbers, proves that the obvious quantity to monitor --- message rate --- is
the \emph{wrong} one to rely on alone, and states the identity that makes an
on-board estimator comparable with a ground-station one.

The motivation is a failure mode already recorded twice on this platform
(Sections~\ref{sec:ovstarvation}, \ref{sec:sqrtvinsint}) and never announced by
any component: when \texttt{/firmware/wheel\_states} slows, MINS loses the
wheel anchoring that gives it its accuracy, and \emph{continues to publish
poses}. The estimator does not fail. It degrades silently into a
visual--inertial estimator with no odometric constraint, which is precisely
the configuration shown elsewhere in this report to diverge. Nothing in the
pipeline distinguishes the two regimes.

\begin{figure}[htbp]
\centering
\begin{tikzpicture}[font=\small, >=Latex, scale=0.95]
  \tikzset{
    blk/.style={draw, rounded corners=1pt, minimum height=8mm,
                align=center, font=\small},
    ok/.style={blk, fill=accentblue!8, draw=accentblue!55},
    bad/.style={blk, fill=alertred!8, draw=alertred!60},
    note/.style={font=\scriptsize\itshape, align=left}
  }
  \node[bad, minimum width=34mm] (cpu)  at (0,0)    {Pi CPU starvation\\
                                                     \scriptsize(SCHED, thermal, contention)};
  \node[bad, minimum width=34mm] (ser)  at (0,-1.6) {\texttt{serial\_node} preempted};
  \node[bad, minimum width=34mm] (whl)  at (0,-3.2) {\texttt{/firmware/wheel\_states}\\
                                                     \scriptsize $18.4 \to 1$\,Hz};
  \node[bad, minimum width=34mm] (mins) at (0,-4.8) {MINS loses wheel anchor};
  \node[ok,  minimum width=34mm] (pub)  at (0,-6.4) {MINS \textbf{still publishes}\\
                                                     \scriptsize poses, covariances, TF};

  \draw[->, thick] (cpu) -- (ser);
  \draw[->, thick] (ser) -- (whl);
  \draw[->, thick] (whl) -- (mins);
  \draw[->, thick] (mins) -- (pub);

  \node[note, anchor=west, color=alertred!85] at (2.6,-3.2)
        {measured 2026-07-28:\\ $18.4 \to 1$\,Hz, UART FIFO overruns};
  \node[note, anchor=west, color=alertred!85] at (2.6,-4.8)
        {Section~\ref{ssec:parallaxe}: without wheels,\\
         $\sigma_Z$ grows as $Z^2$ --- unconstrained};
  \node[note, anchor=west] at (2.6,-6.4)
        {\textbf{no error, no warning, no}\\
         \textbf{covariance inflation.} The\\
         consumer cannot tell the\\
         difference.};

  \draw[decorate, decoration={brace, amplitude=5pt}, color=accentblue]
        (-2.45,0.45) -- (-2.45,-6.85);
  \node[note, anchor=east, color=accentblue, align=right] at (-2.7,-3.2)
        {the entire chain is\\ \emph{invisible} from the\\ consumer's side};
\end{tikzpicture}
\caption{The blind spot that motivates external instrumentation. Each arrow is
a mechanism verified on this platform; the terminal state is indistinguishable
from nominal operation for any node downstream of MINS. A monitor placed
outside the chain is the only way to detect it.}
\label{fig:camp:aveugle}
\end{figure}

\subsection{Rate estimation and the arithmetic of a badly chosen threshold}
\label{ssec:camp:cadence}

\subsubsection{The estimator}

Let a topic deliver messages at arrival instants $t_1 < t_2 < \cdots$, and let
$W$ be the width of a sliding window. Write $N$ for the number of arrivals
inside the current window. The rate estimate used by the monitor is

\begin{equation}
  \hat{f} \;=\; \frac{N-1}{t_N - t_1}.
  \label{eq:camp:fhat}
\end{equation}

The numerator is $N-1$ and not $N$ because the span $t_N - t_1$ covers exactly
$N-1$ inter-arrival intervals; using $N$ inflates the estimate by a factor
$N/(N-1)$, which for a short window is not negligible --- at $N=6$ it is a
$20\,\%$ bias, in the direction that \emph{hides} a slowdown.

\subsubsection{Its dispersion}

Let the inter-arrival intervals $\Delta_k = t_{k+1}-t_k$ be independent with
mean $T = 1/f$ and standard deviation $\sigma_\Delta$, and define the
coefficient of variation $c = \sigma_\Delta / T$, a dimensionless measure of
publication jitter. The span is a sum of $N-1$ such intervals:

\begin{equation}
  S = t_N - t_1 = \sum_{k=1}^{N-1}\Delta_k,
  \qquad
  \mathbb{E}[S] = \frac{N-1}{f},
  \qquad
  \operatorname{Var}(S) = (N-1)\,\sigma_\Delta^2 .
\end{equation}

Since $\hat f = (N-1)/S$, the delta method applied to $g(S)=(N-1)/S$ with
$g'(S) = -(N-1)/S^2$ gives

\begin{equation}
  \operatorname{Var}(\hat f)
  \;\approx\; \left(\frac{N-1}{\mathbb{E}[S]^2}\right)^{\!2}\operatorname{Var}(S)
  \;=\; \frac{f^4}{(N-1)^2}\,(N-1)\,\sigma_\Delta^2
  \;=\; \frac{f^4 \sigma_\Delta^2}{N-1},
\end{equation}

so that the \emph{relative} dispersion takes the compact form

\begin{equation}
  \boxed{\;
  \frac{\sigma_{\hat f}}{f} \;=\; \frac{c}{\sqrt{N-1}}
  \;=\; \frac{c}{\sqrt{fW-1}} \;}
  \label{eq:camp:disp}
\end{equation}

The estimator becomes sharper both with a longer window and with a faster
topic, and the window may therefore be chosen per-topic to hold the dispersion
constant. Table~\ref{tab:camp:disp} evaluates \eqref{eq:camp:disp} at
$W = 5$\,s and $c = 0.10$, a jitter figure consistent with what was observed
on this platform.

\begin{table}[htbp]
\centering
\caption{Dispersion of the rate estimator \eqref{eq:camp:fhat} at $W=5$\,s,
$c=0.10$, and the resulting separation of the $70\,\%$ alarm threshold
expressed in standard deviations.}
\label{tab:camp:disp}
\begin{tabular}{lrrrr}
\toprule
topic & $f$ (Hz) & $N-1$ & $\sigma_{\hat f}/f$ & threshold distance \\
\midrule
\texttt{/firmware/wheel\_states} & 18.4 & 91  & $1.05\,\%$ & $28.6\,\sigma$ \\
\texttt{/imu/data\_clean}        & 83.9 & 419 & $0.49\,\%$ & $61.4\,\sigma$ \\
\texttt{/pc/camera/infra1/\dots} & 15.0 & 74  & $1.16\,\%$ & $25.8\,\sigma$ \\
\bottomrule
\end{tabular}
\end{table}

\subsubsection{Why $20$\,Hz was the wrong threshold, quantitatively}

The nominal rate of the wheel topic is universally cited as $20$\,Hz, and the
initial specification for this monitor was to warn below that figure. The
measured rate, over $30$\,s with the full stack running, is
$f = 18.38$\,Hz.

Treating $\hat f$ as approximately Gaussian --- justified by
$N-1 = 91$ --- with $\sigma_{\hat f} = 0.0105 \times 18.38 = 0.193$\,Hz, the
probability that a healthy stream triggers a threshold placed at
$f_{\mathrm{th}} = 20$\,Hz is

\begin{equation}
  \mathbb{P}\!\left[\hat f < f_{\mathrm{th}}\right]
  = \Phi\!\left(\frac{20 - 18.38}{0.193}\right)
  = \Phi(8.39)
  = 1 - 3\times10^{-17}.
  \label{eq:camp:faux}
\end{equation}

The alarm is not merely frequent: it is certain, at every evaluation, forever.
A monitor in that state is worse than no monitor, because operators learn to
ignore it and it is then silent on the day it matters. Anchoring the threshold
to a \emph{measured} nominal $f_0$ instead, at a fraction $\alpha = 0.70$,
inverts the situation: from \eqref{eq:camp:disp} the separation is

\begin{equation}
  \frac{f_0 - \alpha f_0}{\sigma_{\hat f}}
  = \frac{(1-\alpha)\sqrt{N-1}}{c}
  = \frac{0.30 \times \sqrt{91}}{0.10}
  = 28.6\,\sigma,
\end{equation}

a false-alarm probability below $10^{-179}$. The design rule that follows is
worth stating plainly, because it is the one that was nearly violated:
\emph{a threshold must be derived from the distribution of the quantity it
gates, never from the round number in the specification.}

\subsection{Latency is the leading indicator; rate is the lagging one}
\label{ssec:camp:latence}

\subsubsection{The queueing argument}

Model a consuming node as a single server of service rate $\mu$ fed by arrivals
of rate $\lambda$, with a bounded input queue of $K$ slots --- in ROS, exactly
the subscriber's \texttt{queue\_size}. Write $\rho = \lambda/\mu$ for the
utilisation. Two quantities behave in completely different ways as
$\rho \to 1$.

The mean sojourn time of an unbounded such queue is the classical

\begin{equation}
  \mathbb{E}[W] \;=\; \frac{1}{\mu - \lambda} \;=\; \frac{1/\mu}{1-\rho},
  \label{eq:camp:wait}
\end{equation}

which \emph{diverges} as $\rho\to1^-$; with a bounded queue it saturates at
$K/\mu$, the time to drain a full buffer. The blocking probability of the
bounded queue, on the other hand, is

\begin{equation}
  P_K \;=\; \frac{(1-\rho)\,\rho^{K}}{1-\rho^{K+1}},
  \qquad
  \lim_{\rho\to1} P_K = \frac{1}{K+1},
  \label{eq:camp:pk}
\end{equation}

and message loss --- hence the drop in observed rate --- is exactly $P_K$.
The limit in \eqref{eq:camp:pk} is the crux. It gives a bound that holds over
the \emph{entire} stable regime, and it is worth isolating.

\begin{equation}
  \boxed{\;
  \rho \le 1
  \;\;\Longrightarrow\;\;
  \text{rate loss} \;=\; P_K \;\le\; \frac{1}{K+1}
  \;=\; \frac{1}{51} \;=\; 1.96\,\% \;}
  \label{eq:camp:borne}
\end{equation}

\begin{tcolorbox}[colback=caseboxbg, colframe=accentblue!60,
                  title={Consequence: the rate alarm cannot fire while the
                         queue is stable}]
With $K=50$ and a threshold at $70\,\%$ of nominal, a $30\,\%$ loss is required
to raise an alarm. By \eqref{eq:camp:borne} no utilisation $\rho\le1$ can
produce more than $1.96\,\%$. The rate test is therefore \emph{structurally}
blind to every degree of saturation short of genuine overload. Once
$\rho>1$ the server saturates at throughput $\mu$, the observed rate falls to
$\mu = \lambda/\rho$, and the loss is $1-1/\rho$; requiring that to exceed
$30\,\%$ gives
\[
  1 - \frac{1}{\rho} \ge 0.30
  \quad\Longleftrightarrow\quad
  \rho \ge 1.43 .
\]
The producer must be running $43\,\%$ beyond capacity before the rate test
reacts at all.
\end{tcolorbox}

Latency reaches its threshold far earlier. For the wheel topic
$1/\mu = 1/18.4 = 54.3$\,ms, and the monitor's staleness limit is
$W_{\max} = 0.5$\,s; inverting \eqref{eq:camp:wait},

\begin{equation}
  \frac{1/\mu}{1-\rho} \ge W_{\max}
  \quad\Longleftrightarrow\quad
  \rho \;\ge\; 1 - \frac{1}{\mu W_{\max}}
  \;=\; 1 - \frac{0.0543}{0.5}
  \;=\; 0.891 .
\end{equation}

\begin{table}[htbp]
\centering
\caption{Utilisation at which each test reacts, for
\texttt{/firmware/wheel\_states} ($f=18.4$\,Hz, $K=50$). The ordering is
structural, not a matter of tuning.}
\label{tab:camp:ordre}
\begin{tabular}{llr}
\toprule
test & quantity & triggers at \\
\midrule
staleness $>0.5$\,s        & $\mathbb{E}[W]$, Eq.~\eqref{eq:camp:wait} & $\rho = 0.891$ \\
rate $<70\,\%$ of nominal  & $P_K$ then $1-1/\rho$, Eq.~\eqref{eq:camp:pk} & $\rho = 1.43$ \\
\bottomrule
\end{tabular}
\end{table}

\begin{figure}[htbp]
\centering
\begin{tikzpicture}[font=\small, >=Latex, xscale=13, yscale=0.80]
  \draw[->, thick] (0.48,0) -- (1.07,0) node[below right, font=\scriptsize]
        {$\rho = \lambda/\mu$};
  \draw[->, thick] (0.5,0) -- (0.5,6.0);
  \node[font=\scriptsize, anchor=west, align=left] at (0.487,6.35)
        {test statistic, normalised by its own alarm threshold};
  \foreach \x/\l in {0.5/0.5, 0.7/0.7, 0.9/0.9, 1.0/1.0} {
    \draw (\x,0) -- (\x,-0.16) node[below, font=\scriptsize] {\l};}
  \draw[dashed, color=black!45] (0.5,1) -- (1.055,1)
        node[right, font=\scriptsize] {alarm};

  \draw[thick, color=alertred]
    plot coordinates {(0.50,0.217) (0.60,0.272) (0.70,0.362) (0.75,0.435)
                      (0.80,0.543) (0.84,0.679) (0.87,0.836) (0.891,1.0)
                      (0.91,1.207) (0.93,1.552) (0.95,2.174) (0.96,2.717)
                      (0.97,3.623) (0.976,4.34) (0.9795,5.10)};
  \node[color=alertred, font=\scriptsize, anchor=east, align=right] at (0.958,5.35)
        {staleness $\mathbb{E}[W]/W_{\max}$\\ \emph{diverges as} $(1-\rho)^{-1}$};

  \draw[thick, color=accentblue]
    plot coordinates {(0.50,0.000) (0.70,0.000) (0.80,0.000) (0.90,0.002)
                      (0.95,0.014) (0.99,0.050) (1.00,0.065)};
  \node[color=accentblue, font=\scriptsize, anchor=west] at (0.505,-0.85)
        {rate loss $P_K/0.30$ (blue, flat on the axis): never exceeds
         $0.065$ for any $\rho\le1$};

  \fill[color=alertred] (0.891,1) circle (0.5pt);
  \draw[thin, color=alertred] (0.891,1) -- (0.862,1.62);
  \node[color=alertred, font=\scriptsize, anchor=east] at (0.864,1.70)
        {$\rho^\star = 0.891$};
\end{tikzpicture}
\caption{Why staleness must be monitored alongside rate. Both statistics are
normalised by their own alarm threshold, so the horizontal dashed line at $1$
is where each fires. Staleness crosses it at $\rho^\star = 0.891$; the rate
statistic remains below $0.065$ across the whole stable regime and cannot
cross it at all --- the queue absorbs the pressure as \emph{delay} long before
it sheds it as \emph{loss}.}
\label{fig:camp:queue}
\end{figure}

\subsubsection{The limit of a fixed staleness threshold}

The derivation above makes $\rho^\star$ depend on the topic through
$1/\mu$: a single absolute limit $W_{\max}$ corresponds to
$\rho^\star = 0.891$ on the $18.4$\,Hz wheel stream but to
$\rho^\star = 1 - (0.0119/0.5) = 0.976$ on the $83.9$\,Hz IMU stream, which is
a considerably later warning. The fixed threshold is retained nonetheless, for
a reason outside the queueing model: the measured staleness includes transport
--- the wireless hop and the republication described in
Section~\ref{ssec:camp:comparabilite} --- and half a second of stale odometry
is unusable for control whatever its cause. A per-topic $\rho^\star$ would be
the refinement to make if the monitor were ever used for diagnosis rather than
for vetoing a run.

\subsection{Load average is not utilisation}
\label{ssec:camp:charge}

The decision to run an on-board estimator hinges on how much CPU the Raspberry
Pi has left. The instinct is to read \texttt{load average}; on this platform
that instinct gives the wrong answer by a factor of two.

Linux maintains the load average as an exponentially weighted moving average,
sampled every $\Delta = 5$\,s, over the number $n(t)$ of tasks that are
\emph{either} runnable \emph{or} in uninterruptible sleep:

\begin{equation}
  \ell(t) \;=\; \ell(t-\Delta)\,e^{-\Delta/\tau}
             \;+\; n(t)\left(1 - e^{-\Delta/\tau}\right),
  \qquad \tau \in \{60,\,300,\,900\}\ \text{s}.
  \label{eq:camp:ewma}
\end{equation}

The set counted by $n(t)$ is the whole difficulty. A task in state
\texttt{D} --- blocked on disk, on USB, on a wireless transfer --- contributes
to $\ell$ while consuming \emph{no} CPU whatsoever. Utilisation, by contrast,
is read from the jiffy counters of \texttt{/proc/stat} over an interval:

\begin{equation}
  U \;=\; 1 - \frac{\Delta_{\text{idle}}}{\Delta_{\text{total}}},
  \qquad
  \text{cores free} \;=\; n_{\mathrm{cpu}}\,(1-U).
  \label{eq:camp:util}
\end{equation}

Measured on the rover on 2026-08-24, with the full stack running and no
on-board estimator:

\begin{equation}
  \ell = 4.18 \ \ \text{on}\ n_{\mathrm{cpu}}=4,
  \qquad
  U = 0.54,
  \qquad
  \text{cores free} = 4 \times 0.46 = 1.84 .
\end{equation}

Read as utilisation, $\ell/n_{\mathrm{cpu}} = 1.045$ says the machine is
saturated and the campaign is impossible. The jiffy counters say $46\,\%$ of
four cores is idle. The gap is the population of blocked tasks,

\begin{equation}
  n_{\mathrm{D}} \;\approx\; \ell - n_{\mathrm{cpu}} U
  \;=\; 4.18 - 4 \times 0.54 \;=\; 2.02 ,
\end{equation}

about two tasks permanently waiting on I/O --- consistent with the RealSense
USB pipeline and SD-card traffic, neither of which competes for CPU with an
estimator. Since sqrtVINS runs single-threaded on this build
($\texttt{num\_opencv\_threads}=1$, Section~\ref{sec:ovstarvation}), it
requires approximately one core against the $1.84$ available. The margin is
real but thin, which is why the campaign script measures
\eqref{eq:camp:util} before every run rather than assuming it, and refuses to
start below $25\,\%$ idle --- one core, the strict minimum with no reserve.

\subsection{Comparability: separating the estimator from the architecture}
\label{ssec:camp:comparabilite}

Deploying sqrtVINS on the robot changes what it consumes. The two ground
estimators subscribe to republished streams that have crossed a wireless link;
the on-board instance reads the sensors directly.

\begin{figure}[htbp]
\centering
\begin{tikzpicture}[font=\small, >=Latex, scale=0.95]
  \tikzset{
    blk/.style={draw, rounded corners=1pt, minimum height=7.5mm,
                align=center, font=\scriptsize},
    sens/.style={blk, fill=codebg, draw=codeframe},
    est/.style={blk, fill=accentblue!8, draw=accentblue!55, minimum width=25mm},
    hop/.style={blk, fill=alertred!10, draw=alertred!60},
    zone/.style={draw, dashed, rounded corners=3pt, color=black!35}
  }
  % robot side
  \node[sens, minimum width=22mm] (d455) at (0,0)      {D455 stereo\\ \texttt{/camera/infra\{1,2\}}};
  \node[sens, minimum width=22mm] (imu)  at (0,-1.5)   {CORE2 IMU\\ \texttt{/imu/data\_clean\_pi}};
  \node[est,  minimum width=26mm] (spi)  at (3.7,-0.75){\textbf{sqrtVINS} (Pi)\\ \texttt{/ov\_srvins/odomimu}};
  \draw[->, thick] (d455) -- (spi);
  \draw[->, thick] (imu)  -- (spi);

  % the hop
  \node[hop, minimum width=20mm] (wifi) at (0,-3.5) {WiFi hop\\ + republication};
  \draw[->, thick, color=alertred] (d455.south west) -- ++(-0.35,0) |- (wifi.west);
  \draw[->, thick, color=alertred] (imu.south west)  -- ++(-0.35,0) |- (wifi.west);

  % PC side
  \node[sens, minimum width=24mm] (pcs) at (0,-5.1) {\texttt{/pc/camera/\dots}\\ \texttt{/imu/data\_clean}};
  \draw[->, thick, color=alertred] (wifi) -- (pcs);

  \node[est, minimum width=26mm] (spc) at (3.7,-4.2) {\textbf{sqrtVINS} (PC)\\ \texttt{/sqrtvins/odomimu}};
  \node[est, minimum width=26mm] (ov)  at (3.7,-5.5) {openVINS (PC)\\ \texttt{/ov\_msckf/odomimu}};
  \node[est, minimum width=26mm] (mn)  at (3.7,-6.8) {MINS (PC)\\ \texttt{/mins/imu/odom}};
  \draw[->, thick] (pcs.east) -- (spc.west);
  \draw[->, thick] (pcs.east) -- (ov.west);
  \draw[->, thick] (pcs.east) -- (mn.west);

  \node[zone, fit=(d455)(imu)(spi), inner sep=4mm] {};
  \node[font=\scriptsize\itshape, color=black!55, anchor=west] at (-1.9,0.95)
        {on-board};
  \node[font=\scriptsize\itshape, color=black!55, anchor=west] at (-1.9,-7.6)
        {ground station};

  % the two differences
  \draw[<->, thick, color=accentblue!80] (spi.east) -- ++(0.55,0)
        |- (spc.east);
  \node[font=\scriptsize, color=accentblue, anchor=west, align=left]
        at (5.55,-2.5) {\textbf{same code,}\\ \textbf{different inputs}\\
                        $\Rightarrow$ architecture};
  \draw[<->, thick, color=black!60] (5.30,-4.9) -- (5.30,-6.6);
  \node[font=\scriptsize, anchor=west, align=left] at (5.55,-5.75)
        {\textbf{same inputs,}\\ \textbf{different code}\\
         $\Rightarrow$ estimator};
\end{tikzpicture}
\caption{The two data paths, and why the ground instance of sqrtVINS is
retained even though it is redundant as an estimator. Holding one of the two
factors fixed at a time is what makes the decomposition
\eqref{eq:camp:decomp} available.}
\label{fig:camp:chemins}
\end{figure}

Let $\mathbf{p}_{E,A}(t)$ denote the trajectory produced by estimator $E$ on
architecture $A$, with $E \in \{\text{sqrt},\text{ov},\text{mins}\}$ and
$A \in \{\text{Pi},\text{PC}\}$. Comparing the on-board square-root filter
against the ground openVINS gives a difference in which two independent causes
are superposed. Inserting the ground instance of sqrtVINS --- the same code on
the other architecture --- splits them by a telescoping sum:

\begin{equation}
  \boxed{\;
  \underbrace{\mathbf{p}_{\text{sqrt},\text{Pi}} - \mathbf{p}_{\text{ov},\text{PC}}}
             _{\text{observed}}
  \;=\;
  \underbrace{\mathbf{p}_{\text{sqrt},\text{Pi}} - \mathbf{p}_{\text{sqrt},\text{PC}}}
             _{\text{architecture, code fixed}}
  \;+\;
  \underbrace{\mathbf{p}_{\text{sqrt},\text{PC}} - \mathbf{p}_{\text{ov},\text{PC}}}
             _{\text{estimator, inputs fixed}} \;}
  \label{eq:camp:decomp}
\end{equation}

The identity is trivial; its consequence is not. Without
$\mathbf{p}_{\text{sqrt},\text{PC}}$ the two terms on the right are not
separately observable, and any discrepancy can be attributed to whichever of
the two the reader prefers. This is the reason the campaign keeps
\emph{both} instances of sqrtVINS running and records them as distinct
sources, and the reason the trajectory resampler carries four keys rather than
three. It is also a caution about the campaign's headline comparison: an
on-board estimator measured against ground-based ones is not an
estimator comparison unless the middle term is available.

\subsection{Resampling onto a common time base}
\label{ssec:camp:resample}

The four series arrive on independent clocks at unequal rates. Comparison
requires a common grid, which requires interpolation, which for the rotational
part is not a vector-space operation.

\subsubsection{Rotational interpolation, and the error that actually matters}

Between two unit quaternions $q_i$, $q_{i+1}$ separated by
$\Omega = \arccos\langle q_i, q_{i+1}\rangle$, spherical linear interpolation
is

\begin{equation}
  \operatorname{slerp}(q_i,q_{i+1};u)
  \;=\; \frac{\sin\!\big((1-u)\Omega\big)}{\sin\Omega}\,q_i
      \;+\; \frac{\sin(u\Omega)}{\sin\Omega}\,q_{i+1},
  \qquad u\in[0,1],
  \label{eq:camp:slerp}
\end{equation}

which traverses the geodesic of $S^3$ at constant angular velocity. Normalised
linear interpolation follows the same path but at a non-uniform rate; writing
$\theta(u) = \operatorname{atan2}\!\big(u\sin\Omega,\;(1-u)+u\cos\Omega\big)$
for the angle it actually reaches, a third-order expansion in $\Omega$ gives

\begin{equation}
  \theta(u) - u\Omega
  \;=\; \Omega^3\left(\frac{u^2}{2} - \frac{u}{6} - \frac{u^3}{3}\right)
        + \mathcal{O}(\Omega^5),
  \qquad
  \max_u \big|\theta(u)-u\Omega\big| \approx 0.0161\,\Omega^3 .
  \label{eq:camp:nlerp}
\end{equation}

At $\Omega = \pi/2$ this is $0.062$\,rad, or $3.6^\circ$ --- the familiar
figure that motivates slerp. At the rates in play here it is nothing at all:
the rover's yaw rate stays below $0.5$\,rad\,s$^{-1}$ and odometry is published
at $84$\,Hz, so $\Omega \le 6\times10^{-3}$\,rad and the bound in
\eqref{eq:camp:nlerp} evaluates to $3.5\times10^{-9}$\,rad.

The consequential subtlety is elsewhere. The unit quaternions double-cover
$SO(3)$: $q$ and $-q$ denote the \emph{same} rotation, and a publisher may emit
either. If consecutive samples straddle a sign change,
$\langle q_i,q_{i+1}\rangle < 0$, then \eqref{eq:camp:slerp} interpolates along
the long arc, sweeping $2\pi - \Omega$ instead of $\Omega$. The remedy is one
line --- replace $q_{i+1}$ by $-q_{i+1}$ whenever the inner product is
negative --- and the asymmetry of the two errors is the point:

\begin{equation}
  \frac{\text{sign-flip error}}{\text{slerp-versus-lerp error}}
  \;\approx\; \frac{\pi}{3.5\times10^{-9}}
  \;\approx\; 9\times10^{8}.
\end{equation}

Nine orders of magnitude separate the failure everyone guards against from the
one that is easy to omit. The validation written for this pipeline initially
tested $q_w \approx 1$ at the origin, which a legitimately negated quaternion
fails; the correct invariant is $|q_w| \approx 1$.

\subsubsection{Origin alignment, and why its order is not free}

Each series is expressed in its own initialisation frame. Alignment to a common
origin is the rigid transformation

\begin{equation}
  \mathbf{p}'(t) = R_0^{\top}\big(\mathbf{p}(t) - \mathbf{p}_0\big),
  \qquad
  q'(t) = q_0^{-1} \otimes q(t),
  \label{eq:camp:se3}
\end{equation}

with $(\mathbf{p}_0, q_0)$ the pose at the origin of the comparison. The
operation must be applied \emph{after} resampling, not before, and the reason
is a timing argument rather than an algebraic one.

Let series $X$ have raw samples beginning at $t^X_0$, and let the common window
begin at $T = \max_X t^X_0$, since the intersection of the supports --- not
their union --- is the only interval on which every series is a measurement
rather than an extrapolation. Aligning before resampling sets
$(\mathbf{p}_0,q_0) = (\mathbf{p}_X(t^X_0), q_X(t^X_0))$, so that on the common
grid

\begin{equation}
  \mathbf{p}'_X(T) = R_0^{\top}\big(\mathbf{p}_X(T) - \mathbf{p}_X(t^X_0)\big)
  \;\ne\; \mathbf{0}
  \qquad\text{whenever } t^X_0 < T,
\end{equation}

which is every series except the latest to start. The trajectories then begin
at different points, and since the campaign's primary metric is loop closure
--- the norm $\|\mathbf{p}(t_{\text{end}})-\mathbf{p}(t_{\text{start}})\|$ for
a path returning to its start --- an origin offset enters the result directly
and silently. Applying \eqref{eq:camp:se3} after resampling, with
$(\mathbf{p}_0,q_0)$ taken at $T$, makes every series start at the origin by
construction.

\subsection{The campaign protocol}
\label{ssec:camp:protocole}

\subsubsection{Ordering, and the coverage condition}

The launch order is not a convenience. Write $[s,e]$ for the temporal support
of the recording and $[m_s,m_e]$ for that of the health monitor. The recording
is admissible as evidence only if

\begin{equation}
  m_s \le s \quad\text{and}\quad m_e \ge e,
  \label{eq:camp:couverture}
\end{equation}

that is, only if every instant of the bag is covered by health evidence.
Violating the left inequality is the realistic failure: starting the monitor
after the recorder leaves an initial interval --- exactly when nodes are
starting and contention peaks --- with no observation, and a starvation
occurring there would corrupt a file believed to be clean. Since the monitor
additionally needs one window $W$ to fill before it can judge, the operational
requirement is $m_s \le s - W$.

\begin{figure}[htbp]
\centering
\begin{tikzpicture}[font=\small, >=Latex, xscale=1.24, yscale=1.0]
  \tikzset{
    etape/.style={draw, rounded corners=1pt, minimum height=7mm,
                 align=center, font=\scriptsize, fill=codebg, draw=codeframe},
    gate/.style={etape, fill=alertred!8, draw=alertred!60}
  }
  \node[gate,  minimum width=19mm] (v1) at (0.95,2.4) {\textbf{1} ground checks\\
                                                       \scriptsize temp, CPU, battery};
  \node[gate,  minimum width=19mm] (v2) at (2.95,2.4) {\textbf{1bis} robot checks\\
                                                       \scriptsize idle, throttle, disk};
  \node[etape, minimum width=18mm] (s3) at (4.90,2.4) {\textbf{2} sqrtVINS (Pi)\\
                                                       \scriptsize slowest to start};
  \node[etape, minimum width=18mm] (s4) at (6.75,2.4) {\textbf{3} monitor\\
                                                       \scriptsize + Pi telemetry};
  \node[etape, minimum width=15mm] (s5) at (8.50,2.4) {\textbf{4} \texttt{rosbag}};

  \foreach \a/\b in {v1/v2, v2/s3, s3/s4, s4/s5}
    \draw[->, thick] (\a) -- (\b);

  % support du moniteur
  \draw[thick, color=accentblue] (6.75,1.30) -- (9.30,1.30);
  \draw[thick, color=accentblue] (6.75,1.16) -- (6.75,1.44);
  \draw[thick, color=accentblue] (9.30,1.16) -- (9.30,1.44);
  \node[font=\scriptsize, color=accentblue, anchor=east] at (6.58,1.30)
        {$[m_s,m_e]$ monitor};

  % avance requise
  \draw[<->, color=black!60] (6.78,0.78) -- (7.92,0.78);
  \node[font=\scriptsize, color=black!60, anchor=south] at (7.35,0.84) {$\ge W$};

  % support de l'enregistrement
  \draw[thick, color=alertred] (7.95,0.28) -- (9.10,0.28);
  \draw[thick, color=alertred] (7.95,0.14) -- (7.95,0.42);
  \draw[thick, color=alertred] (9.10,0.14) -- (9.10,0.42);
  \node[font=\scriptsize, color=alertred, anchor=east] at (7.78,0.28)
        {$[s,e]$ bag};

  \draw[->, thick] (0.0,-0.35) -- (9.7,-0.35)
        node[below right, font=\scriptsize] {$t$};

  \node[font=\scriptsize\itshape, anchor=west, align=left] at (0.0,-1.05)
        {a run whose checks fail is \emph{refused}, not flagged: an inadmissible\\
         measurement costs more to interpret than to repeat};
\end{tikzpicture}
\caption{Campaign sequence. Stages 1 and 1bis are vetoes. The monitor's support
strictly contains the recording's, with at least one window $W$ of lead so that
its sliding estimators are converged before the first message is written ---
condition \eqref{eq:camp:couverture}.}
\label{fig:camp:sequence}
\end{figure}

\subsubsection{Admission thresholds}

Table~\ref{tab:camp:seuils} collects the conditions under which the campaign
script refuses to start, each with the mechanism that justifies it. Two of them
--- ground temperature and battery voltage --- were already established in
Section~\ref{ssec:bilan} after they invalidated earlier runs; the three robot
conditions are added here.

\begin{table}[htbp]
\centering
\caption{Admission conditions. A run failing any of them is refused outright.}
\label{tab:camp:seuils}
\small
\begin{tabular}{p{3.1cm} p{1.9cm} p{7.6cm}}
\toprule
quantity & veto & mechanism \\
\midrule
ground CPU temperature & $\ge 82^{\circ}$C & thermal throttling; the recorder
falls behind and the loss is written into the bag indistinguishably from a
sensor fault \\[2pt]
browser CPU            & $\ge 50\,\%$      & measured cause of the above; the
cockpit overlay alone reached $109\,\%$ \\[2pt]
battery                & $<11.1$\,V        & below $\approx30\,\%$ charge yaw
authority degrades, so a commanded rotation is not the rotation performed \\[2pt]
robot idle             & $<25\,\%$         & $25\,\%$ of four cores is one core;
sqrtVINS is single-threaded and needs it, Eq.~\eqref{eq:camp:util} \\[2pt]
robot \texttt{throttled} & $\ne$ \texttt{0x0} & the SoC has clocked down; this
is the mechanism behind the UART FIFO overruns of 2026-07-28 \\[2pt]
robot disk             & $\ge 92\,\%$      & a full card has already crashed
\texttt{leo.service} on this platform \\
\bottomrule
\end{tabular}
\end{table}

\subsubsection{Where each observation is made}

The monitor runs on the ground station and not on the robot. The reasoning is
that a monitor built to detect CPU starvation must not contribute to the load
it measures: subscribing to an $84$\,Hz stream costs real CPU on the very
machine under test. All topics are visible from the ground through the shared
master, so nothing is lost by observing from outside.

That argument does not extend to the robot's own resource counters, and the
distinction is worth making explicit because it looks like an exception.
Reading \texttt{/proc/stat} and the SoC's throttling register every two seconds
costs no measurable CPU --- there is no subscription, no deserialisation, no
callback. It is also the only place the \emph{cause} can be seen: the ground
monitor observes effects, a rate that falls, while saturation and clock
throttling are visible only on the machine that suffers them. The campaign
therefore samples both, and the post-run summary reports the minimum idle
fraction alongside the rate transitions so the two can be correlated.

The final verification closes the loop on the recording itself. Rates are
recomputed \emph{from the bag} rather than taken from any live counter: a topic
can publish at $18.4$\,Hz and still be written at $12$, if the recorder was the
component under pressure. Only the file's own contents establish what was
captured, which is the whole purpose of an admissibility check ---
Eq.~\eqref{eq:camp:fhat} applied to the recording, with
$N$ the message count and $t_N - t_1$ the bag duration.
